\chapter{Lições Aprendidas}

O projeto Micromouse, da disciplina Projeto Integrador 1 (PI1), proporcionou  valiosas lições
para o desenvolvimento de futuros projetos interdisciplinares. A integração entre diversas áreas
das Engenharias de Software, Aeroespacial, Eletrônica, Energia e Automotiva, demonstrou o potencial da 
colaboração multidisciplinar.

Este capítulo documenta as lições aprendidas pela equipe ao longo do projeto, abordando desde a análise do cenário até os desafios de execução, visando evitar a repetição de falhas e otimizar processos em projetos futuros.

\section{Análise SWOT}

\begin{table}[ht]
% \label{tab:swot}
\centering
\caption{Análise SWOT do produto.}
\begin{tabular}{|c|c|c|} \hline
 & \textbf{Pontos fortes} & \textbf{Pontos fracos}  \\ \hline 
    {\parbox{0.1\textwidth}{ \textbf{Fatores \\ internos}}}  &
    {\parbox{0.4\textwidth}{ Forças: 
        \begin{itemize}
            \item Respeito e colaboração interdisciplinar, com subgrupos se ajudando ativamente.
            \item Resiliência da equipe diante da necessidade de troca de componentes.
            \item Decisão de projeto robusta com a utilização de chassi em dois andares.
        \end{itemize} }} & 
    {\parbox{0.4\textwidth}{ Fraquezas:
        \begin{itemize}
            \item Falha de comunicação e dificuldade na integração entre as subequipes.
            \item Desorganização no cumprimento de prazos, sobretudo na entrega do chassi.
            \item Desequilíbrio na distribuição de tarefas entre os integrantes do grupo.
        \end{itemize} }}\\ \hline
    {\parbox{0.1\textwidth}{ \textbf{Fatores \\ externos}}}  &
    {\parbox{0.4\textwidth}{ Oportunidades:
        \begin{itemize}
            \item Aprimoramento de \textit{soft skills} vinculadas à gestão de projetos.
            \item Aprendizado técnico prático em áreas até então desconhecidas pelos membros.
            \item Enriquecimento sistêmico através de ambiente multidisciplinar entre engenharias.
        \end{itemize} }} & 
    {\parbox{0.4\textwidth}{ Ameaças:
        \begin{itemize}
            \item Dificuldade e atrasos no fornecimento de recursos estruturais para o robô.
            \item Competição de tempo dos alunos devido aos compromissos com outras disciplinas.
            \item Riscos de falha no projeto causados pela sobrecarga de alguns integrantes.
        \end{itemize} }}\\ \hline
\end{tabular}
\label{tab:swot}
\end{table}


\section{Planejado x realizado}
\subsection{Objetivos}
Os objetivos do projeto foram plenamente alcançados. O micromouse demonstrou autonomia e eficácia na resolução do labirinto durante a apresentação 
final, validando a aplicação prática do algoritmo de navegação. A integração física e lógica entre os subsistemas de estrutura, eletrônica e firmware
funcionou de maneira harmônica, garantindo leituras estáveis dos sensores e respostas dos motores em tempo real. Essa conclusão com sucesso consolidou
o aprendizado prático da equipe e superou os gargalos de comunicação e atrasos enfrentados nas etapas iniciais de desenvolvimento.

\subsection{Prazos}

Ocorreram atrasos pontuais, sobretudo pela alta dependência entre as equipes e na entrega do chassi final utilizado. Isso gerou uma concentração do trabalho
do subgrupo de hardware e software na reta final do semestre em poucos integrantes.

\subsection{Orçamento}

O projeto foi executado com um custo financeiro menor que o inicialmente estipulado. Essa economia foi possível graças ao esforço da equipe no reaproveitamento de materiais
já existentes e no uso de componentes reciclados para a estrutura do labirinto utilizado nos testes da equipe. No entanto, houve a necessidade de troca de bateria algumas
vezes no decorrer do semestre.

\subsection{Escopo}

O projeto atendeu o escopo definido pelos requisitos da disciplina. Todas as funcionalidades e requisitos foram alcançados e o micromouse passou na avaliação final
(resolver dois labirintos 4x4 e um 8x8), obtendo a aprovação dos professores. Apesar de alguns atrasos em tarefas dentro do projeto, ele respeitou o tempo da entrega
final definido no cronograma da disciplina.


\section{Projeto}
\subsection{Pontos fortes}
\begin{enumerate}
    \item A liberdade de escolha dos membros em qual equipe atuar, facilitou a execução e resolução das atividades devido a familiaridade com a área.
    \item A impressão de um protótipo para testes permitiu que o grupo pudesse realizar ajustes e otimizações estruturais do micromouse.
    \item O uso de ferramentas já conhecidas (GitHub, LaTeX) por boa parte da equipe evitou uma grande curva de aprendizado. 
    \item Fazer múltiplos testes em todas as aulas e comunicar o grupo sobre o andamento nos deixou mais preparados para a apresentação.
\end{enumerate}

\subsection{Pontos fracos}
\begin{enumerate}
    \item Comunicação ineficiente e falta de fluxo de informações contínuo entre as subequipes.
    \item Ausência de uma gerência rigorosa e centralizada para a cobrança e acompanhamento de prazos.
    \item Desbalanceamento na distribuição de tarefas gerou sobrecarga em alguns membros na reta final.
    \item A falta de integração trouxe atraso em algumas atividades que possuem dependência entre diferentes subequipes.
\end{enumerate}

\section{Recomendações para projetos futuros}

\begin{enumerate}
    \item Liderança mais ativa no acompanhamento e mais rígida em relação aos prazos estipulados.
    \item A equipe deve melhorar a comunicação de problemas e dar feedbacks constantes das atividades.
    \item Aproveitar melhor ferramentas de gerenciamento de equipe para evitar desbalanceamento.
    \item Incentivar maior sinergia e reuniões periódicas entre equipes para melhorar o planejamento das atividades.
\end{enumerate}

% \section{Questões em aberto}

% \textcolor{red}{Usar caso haja alguma questão pendente em relação às entregas do projeto (Ex.: Requisitos não entregues, melhoria na programação do software, troca de sensores, entre outros.)}

% \begin{enumerate}
%     \item \textcolor{red}{Até 200 caracteres, considerando os espaços.} 
%     \item \textcolor{red}{Até 200 caracteres, considerando os espaços.} 
%     \item \textcolor{red}{Até 200 caracteres, considerando os espaços.} 
%     \item \textcolor{red}{Até 200 caracteres, considerando os espaços.} 
% \end{enumerate}

\section{Desempenho dos fornecedores}

\subsection{Fornecedores com desempenho acima do esperado}

A VoluMaker (@volumaker3d) entregou com rapidez e qualidade a impressão 3D do primeiro protótipo do chassi, peça essencial para testes de posição dos componentes e otimizações estruturais.

\subsection{Fornecedores com desempenho conforme esperado}

Realizamos a compra de diversos componentes eletrônicos por meio da loja HU Infinito, sendo que todos os produtos apresentaram a qualidade esperada e foram entregues dentro do prazo estipulado.

\subsection{Fornecedores com desempenho abaixo do esperado}

Tivemos problema na compra de um kit de parafusos pelo Mercado Livre, pois a entrega não foi realizada, atrasando a montagem e acarretando na necessidade de solicitação de reembolso.
